\section{FSA-v3: Bimodal Training Extension} \label{sec:fsa-v3}

The FSA-v3 extension generalizes the model to concurrent training regimes by introducing a fourth latent state representing Strength/Resistance adaptation ($S$) and a second control input ($\Phi_S$).

\subsection{4D State Space Dynamics}

The state vector is expanded to $y(t) = (B, S, F, A)^T$. The dynamics are governed by the following system of coupled SDEs:

\begin{align}
  \dd B &= \left[\kappa_B \cdot \Gamma_B(A) \Phi_B(t) - \frac{B}{\tau_B}\right] \dd t + \sigma_B \sqrt{B(1-B)} \dd W_B \\
  \dd S &= \left[\kappa_S \cdot \Gamma_S(A) \Phi_S(t) - \frac{S}{\tau_S}\right] \dd t + \sigma_S \sqrt{S(1-S)} \dd W_S \\
  \dd F &= \left[\kappa_{F,B} \Phi_B(t) + \kappa_{F,S} \Phi_S(t) - \frac{\Lambda(A)}{\tau_F} F\right] \dd t + \sigma_F \sqrt{F} \dd W_F \\
  \dd A &= \left[\mu(B, S, F) A - \eta A^3\right] \dd t + \sigma_A \sqrt{A} \dd W_A
\end{align}

where the bimodal bifurcation parameter $\mu$ is defined as:
\begin{equation}
  \mu(B, S, F) = \mu_0 + \mu_B B + \mu_S S - \mu_F F - \mu_{FF}(F - F_{\mathrm{typ}})^2
\end{equation}

And the autonomic coupling factors are:
\begin{equation}
    \Gamma_i(A) = \frac{1 + \epsilon_{A,i} A}{1 + \epsilon_{A,i} A_{\mathrm{typ}}}, \quad \Lambda(A) = \frac{1 + \lambda_A A}{1 + \lambda_A A_{\mathrm{typ}}}
\end{equation}

\subsection{Unified Fatigue and Autonomic Robustness}

A key architectural feature of FSA-v3 is the **Unified Fatigue Pool** ($F$). Both aerobic ($\Phi_B$) and strength ($\Phi_S$) stimuli contribute to the same systemic fatigue state, albeit with different gains ($\kappa_{F,B}, \kappa_{F,S}$). This allows the high-frequency aerobic observations (HR, Sleep) to indirectly inform the fatigue state generated by sparse strength sessions.

Furthermore, strength adaptation $S$ is coupled to the autonomic growth rate via $\mu_S$, reflecting the chronotropic effect of neuromuscular robustness on autonomic stability.

\subsection{Bimodal Observation Model}

To identify the strength compartment, a fifth observation channel is introduced:
\begin{equation}
    \mathrm{VolumeLoad}(t) \sim \mathcal{N}(\beta_S S(t) - \beta_F^{VL} F(t), \sigma_{VL}^2)
\end{equation}
Volume Load (Sets $\times$ Reps $\times$ Intensity) provides the primary signal for $S$, while its dependence on $F$ captures the acute performance suppression typical of high-fatigue states.
